

\documentclass{article}
\usepackage[utf8]{inputenc}
% Better font encoding for French and scalable fonts
\usepackage[T1]{fontenc}
\usepackage{lmodern}
\usepackage{amsmath, amssymb, amsfonts, amsthm, stmaryrd}
\usepackage{graphicx}
\usepackage{booktabs}
\usepackage{geometry}
\usepackage[french]{babel}
\usepackage{amsopn}

\title{Etat de l'art des méthodes d'analyse des RSO par apprentissage supervisé et non supervisé}
\author{Etienne Chevrolat}
\date{}


\begin{document}

\maketitle

\begin{abstract}
    Dans ce document on cherche à documenter et résumer l'ensemble des méthodes de Deep learning sur Time-Series explorées lors de mon stage chez Aldoria.
    Mon but principal était de détecter des manoeuvres à partir de séries temporelles de TLEs. 
    Après une analyse des méthodes statistiques existantes, 
    la question de trouver une méthode plus générale/plus fiable/ne necéssitant pas d'analyse au cas par cas/ne nécessitant pas de choix d'un seuil de bruit/anomalies fait naturellement s'intéresser à des méthodes d'apprentissage statistique.  
    On s'est alors plus généralement intéressé à l'analyse globale du catalogue RSO : détection d'anomalies, caractérisation de patterns-of-life, embedding des objets, clusterisation des objets en orbite basse, reconstruction et forecasting des séries temporelles d'éléments orbitaux.
\end{abstract}

\section{Apprentissage supervisé}

\subsection{La difficulté de trouver des données annotées pour notre usage}
    Le principal défi lorsqu'on veut faire de la détection de manoeuvres sur des séries temporelles de paramètres orbitaux par apprentissage, c'est le manque de données suffisantes des manoeuvres de satellites.
    La principale source de données publiques pour les manoeuvres de satellites en orbite basse sont celles des satellites équipés de sondes DORIS : problème, il s'agit d'une poignée de satellites scientifiques européens aux comportements uniques et non généralisables à la plupart des satellites.
    J'ai par ailleurs conçu un dataset personnalisé de maneouvres labellisées à la main de 75 objets d'orbites diverses. D'autres sources de données privées précises pour quelques satellites ont été trouvées.
    Dans tous les cas, la quantité et/ou la qualité des données était problématique. 

\subsection{La performance de petits modèles sur données synthétiques}
    Une étude approfondie de la performance sur des données synthétiques SPLID (osculatrices, non-moyennées) d'objets en orbite GEO à été menée.
    On a ainsi réussi à atteindre les perfs du gagnant du concours SPLID sur les taches de localisation des points de changement de pattern of life + classification des noeuds (début de drift, début de station keeping, ajustement de drift, station keeping electrique/chimique/hybride).
    Les architectures qui fonctionnent sont basiques : fenêtrage des données (128 pas passés, 48 pas futur), un réseau de neurones convolutif (CNN) pour localisation et un réseau récurrent type long-short-term-memory (LSTM) pour la classification.
    

    Une étude avait déjà été réalisée chez Aldoria sur des données synthétiques LEO et des très bon résultat avait aussi été obtenu.
    La pertinence de faire du machine learning sur les données orbitales est donc confirmée : que ce soit sur TLEs ou sur SP Vectors, les modèles d'apprentissage/deep learning peuvent apprendre les patterns des satellites (d'autant plus qu'avec le nombre croissant d'objets déployés, les manoeuvres seront prochainement massivement automatisées et donc prédictibles)

    Le problème est toujours le passage de données synthétiques vers les données réelles, notamment à des méthodes qui fonctionnent sur les données SpaceTrack : des TLEs fiables, réelles et disponible librement en grande quantité.

\section{Apprentissage non supervisé}
    On s'est grandement interessé à cause du manque de données annotées à des méthodes d'apprentissage non-supervisé.
    La source de données est les TLES SpaceTrack de l'ensemble des objets en orbite basse de 2024 à aujourd'hui.
    \subsection{Une phase de pretraining sur les données SpaceTrack}
    (cf article Guimareaes)

    L'idée est d'entrainer un encodeur de type VisionTransformer (basé sur le mécanisme d'attention temporelle cf attention is all you need, cf ViT) par la technique de masquage : 
    On s'est inspiré de ce qui se faisait en préentrainement non-supervision en Vision : cf videoMAE, videoMAEv2 

    On divise la série temporelle en fenêtres (de 2 semaines à 3 mois) et les fenêtres en patchs (de quelques jours).
    

    Un certain nombre des patchs (> 50 \% )de chaque fenetre sont masqués, et l'encodeur applique un embedding (transformers sur un espace latent de dimension 64,128) sur les patchs visibles uniquement.
    Un décodeur reconstruit les patchs masqués et la loss MSE est calculée entre les données masquées et les données réelles uniquement (objectif reconstruction).
    Ces méthodes ont initialement été inventées pour diminuer la taille des données à traiter, car les encodeurs ont un temps de calcul en $O(n^ 2)$ avec $n$ la taille de l'entrée
    Il s'avère que cette technique de masking permet à l'encodeur d'apprendre les patterns généraux sous-jacents aux données.

    L'encodeur ainsi préentrainé (un ViT classique) est ainsi utilisable pour un grand nombre de tache : l'idée est qu'il à appris les mécanismes des données orbitales LEO de manière générale : différencier les types d'orbites, différencier les comportements débris/payload, les fréquence des manoeuvres de station keeping, etc

    On réutilise ensuite cet encodeur pour des tâches spé

\end{document}